\documentclass{article}

% Language setting
\usepackage[english]{babel}

% Set page size and margins
\usepackage[letterpaper,top=2cm,bottom=2cm,left=3cm,right=3cm,marginparwidth=1.75cm]{geometry}

% Useful packages
\usepackage{amsmath}
\usepackage{amssymb}
\usepackage{graphicx}
\usepackage[colorlinks=true, allcolors=black]{hyperref}
\usepackage{titlesec}
\usepackage{xparse}

\NewDocumentCommand{\qcq}{m o o m}{%
	\ensuremath{\forall #1 \IfValueT{#2}{, #2} \IfValueT{#3}{, #3}  \in  \mathbb{#4}}%
}
\NewDocumentCommand{\ext}{m o o m}{%
	\ensuremath{\exists #1 \IfValueT{#2}{, #2} \IfValueT{#3}{, #3}  \in  \mathbb{#4}}%
}

\title{\textbf{Exercices : Arithmétique dans $\mathbb{Z}$}}
\author{\textbf{Yassin Akkab}}

\newcommand{\R}{\mathbb{R}}
\newcommand{\N}{\mathbb{N}}
\newcommand{\Z}{\mathbb{Z}}
\newcommand{\C}{\mathbb{C}}
\newcommand{\Q}{\mathbb{Q}}
\newcommand{\D}{\mathbb{D}}

\begin{document}
	
	\maketitle
	
	\subsection*{Exercice 1 (Division euclidienne et divisibilité)}
	\begin{enumerate}
		\item Montrer que pour tout entier relatif $n$, le nombre $n(n+1)(2n+1)$ est divisible par 6.
		\item Déterminer tous les entiers relatifs $n$ tels que $n-3$ divise $n+5$.
	\end{enumerate}
	
	\subsection*{Exercice 2 (Algorithme d'Euclide et Bézout)}
	Soit l'équation diophantienne suivante dans $\mathbb{Z}^2$ :
	\[ (E) : 13x - 5y = 1 \]
	\begin{enumerate}
		\item Justifier que l'équation $(E)$ admet des solutions dans $\mathbb{Z}^2$.
		\item Trouver une solution particulière $(x_0, y_0)$ de $(E)$.
		\item Résoudre l'équation $(E)$ dans $\mathbb{Z}^2$.
	\end{enumerate}
	
	\subsection*{Exercice 3 (Congruences et restes)}
	\begin{enumerate}
		\item Déterminer le reste de la division euclidienne de $3^{2026}$ par 7.
		\item Montrer que pour tout entier naturel $n$, le nombre $3^{2n+1} + 2^{n+2}$ est divisible par 7.
	\end{enumerate}
	
	\subsection*{Exercice 4 (Théorème de Fermat)}
	Soit $p$ un nombre premier impair.
	\begin{enumerate}
		\item Montrer que pour tout entier $x$, $x^p \equiv x \pmod p$.
		\item Soit $a$ et $b$ deux entiers tels que $a^p \equiv b^p \pmod p$. Montrer que $a \equiv b \pmod p$.
	\end{enumerate}
	
\end{document}
